\chapter{反应机理：弯箭头}\label{ch:g12:curly-arrows}

\omterm{def:g8:balanced-equations:equation}{配平的方程式}就像电影的第一张和最后一张照片：
它展示反应前的\omterm{def:g7:chemical-reactions:reactant}{反应物}和反应后的生成物，
而两者之间的过程一概不显示。然而，\omterm{def:g7:atoms-and-molecules:molecule}{分子}并不是简单地交换\omterm{def:g7:atoms-and-molecules:atom}{原子}。
键是一次断裂或形成一个，\omterm{def:g9:inside-the-atom:nucleus}{电子}对在不同位置之间移动，
经过一些寿命很短、从不出现在方程式中的物种。
化学家用\omterm{def:g12:curly-arrows:curly-arrow}{弯箭头}来表示这些移动。
掌握几条简单规则后，这些箭头便能解释反应为何在某处发生，
并预测从未见过的反应。

\begin{recall}
\omterm{def:g11:polarity-and-cohesion:electronegativity}{电负性}、\omterm{def:g11:polarity-and-cohesion:polar-bond}{极性键}和\omterm{def:g11:polarity-and-cohesion:polar-bond}{部分电荷}
（\cref{ch:g11:polarity-and-cohesion}）。\omterm{def:g10:lewis-and-shape:lewis-structure}{路易斯结构}：\omterm{def:g10:lewis-and-shape:lone-pair}{成键电子对}和
\omterm{def:g10:lewis-and-shape:lone-pair}{孤电子对}（\cref{ch:g10:lewis-and-shape}）。\omterm{def:g11:functional-groups:functional-group}{官能团}与有机物
类别（\cref{ch:g11:functional-groups}）。
\end{recall}

\section{富电子位点与缺电子位点}

\begin{definition}[给电子位点与受电子位点]\label{def:g12:curly-arrows:site}
\omterm{def:g7:atoms-and-molecules:molecule}{分子}或\omterm{def:g9:ions:ion}{离子}中的\emph{给电子位点}\index{给电子位点}是富\omterm{def:g9:inside-the-atom:nucleus}{电子}的位置，
能够提供一对\omterm{def:g9:inside-the-atom:nucleus}{电子}以形成新键，包括\omterm{def:g10:lewis-and-shape:lone-pair}{孤电子对}、\omterm{def:g10:lewis-and-shape:multiple-bond}{双键}、
带负电荷的\omterm{def:g7:atoms-and-molecules:atom}{原子}或
带部分负电荷（$\delta^-$）的\omterm{def:g7:atoms-and-molecules:atom}{原子}。\emph{受电子位点}\index{受电子位点}是缺\omterm{def:g9:inside-the-atom:nucleus}{电子}的位置，
能够接受一对\omterm{def:g9:inside-the-atom:nucleus}{电子}，即带正电荷的\omterm{def:g7:atoms-and-molecules:atom}{原子}或
带部分正电荷的\omterm{def:g7:atoms-and-molecules:atom}{原子}
（$\delta^+$）。
\end{definition}

\begin{example}[两个分子中的位点]\label{ex:g12:curly-arrows:sites}
在溴乙烷 \ce{CH3-CH2-Br} 中，溴比碳的\omterm{def:g11:polarity-and-cohesion:electronegativity}{电负性}大，与溴相连的碳带 $\delta^+$（\omterm{def:g12:curly-arrows:site}{受电子位点}），
溴带 $\delta^-$，并有三对\omterm{def:g10:lewis-and-shape:lone-pair}{孤电子对}。
在丙酮中，\ce{C=O} 基团的氧带 $\delta^-$，
有两对\omterm{def:g10:lewis-and-shape:lone-pair}{孤电子对}（\omterm{def:g12:curly-arrows:site}{给电子位点}），其碳带 $\delta^+$（\omterm{def:g12:curly-arrows:site}{受电子位点}）。
氢氧根\omterm{def:g9:ions:ion}{离子} \ce{OH-} 带负电荷、有三对\omterm{def:g10:lewis-and-shape:lone-pair}{孤电子对}，
是较强的给\omterm{def:g9:inside-the-atom:nucleus}{电子}体。
\end{example}

\begin{omfigure}
\setchemfig{atom sep=2.2em}
\begin{tabular}{cc}
\chemfig{H_3C-C(-[2]H)(-[6]H)-\charge{90:2pt=\:,0=\:,270:2pt=\:}{Br}} &
\chemfig{H_3C-C(-[6]CH_3)=[2]\charge{0=\:,180=\:}{O}} \\[2pt]
\small $\delta^+$ 位于与 \ce{Br} 相连的碳上，$\delta^-$ 位于 \ce{Br} 上 &
\small $\delta^+$ 位于 \ce{C=O} 的碳上，$\delta^-$ 位于 \ce{O} 上 \\
\small 溴乙烷 & \small 丙酮
\end{tabular}

{\small 缺\omterm{def:g9:inside-the-atom:nucleus}{电子}的碳\omterm{def:g7:atoms-and-molecules:atom}{原子}（$\delta^+$，\omterm{def:g12:curly-arrows:site}{受电子位点}）与
富\omterm{def:g9:inside-the-atom:nucleus}{电子}、\omterm{def:g11:polarity-and-cohesion:electronegativity}{电负性}大的\omterm{def:g7:atoms-and-molecules:atom}{原子}（$\delta^-$，带\omterm{def:g10:lewis-and-shape:lone-pair}{孤电子对}，
即\omterm{def:g12:curly-arrows:site}{给电子位点}）相邻。}
\end{omfigure}

\section{弯箭头}

\begin{definition}[弯箭头]\label{def:g12:curly-arrows:curly-arrow}
\emph{弯箭头}\index{弯箭头}表示反应某一步中一对\omterm{def:g9:inside-the-atom:nucleus}{电子}的移动，
它始于移动的那一对\omterm{def:g9:inside-the-atom:nucleus}{电子}，即\omterm{def:g10:lewis-and-shape:lone-pair}{孤电子对}或化学键，
终于\omterm{def:g9:inside-the-atom:nucleus}{电子}所去的位置：
既可以是在某\omterm{def:g7:atoms-and-molecules:atom}{原子}上形成新键，
也可以是接受断裂之键的\omterm{def:g9:inside-the-atom:nucleus}{电子}对的\omterm{def:g7:atoms-and-molecules:atom}{原子}。
\end{definition}

\begin{method}[画弯箭头]\label{met:g12:curly-arrows:drawing}
\begin{enumerate}
  \item 找出\omterm{def:g12:curly-arrows:site}{给电子位点}和\omterm{def:g12:curly-arrows:site}{受电子位点}。
  \item 从给\omterm{def:g9:inside-the-atom:nucleus}{电子}的\omterm{def:g9:inside-the-atom:nucleus}{电子}对（\omterm{def:g10:lewis-and-shape:lone-pair}{孤电子对}或化学键）画一个箭头指向
        受\omterm{def:g9:inside-the-atom:nucleus}{电子}\omterm{def:g7:atoms-and-molecules:atom}{原子}：该处形成新键。
  \item 若受\omterm{def:g9:inside-the-atom:nucleus}{电子}\omterm{def:g7:atoms-and-molecules:atom}{原子}此时的\omterm{def:g9:inside-the-atom:nucleus}{电子}数过多（超过八\omterm{def:g9:inside-the-atom:nucleus}{电子}，
        对氢\omterm{def:g7:atoms-and-molecules:atom}{原子}则超过两个\omterm{def:g9:inside-the-atom:nucleus}{电子}），再画一个箭头断开它的一个键，
        \omterm{def:g9:inside-the-atom:nucleus}{电子}对归于\omterm{def:g11:polarity-and-cohesion:electronegativity}{电负性}较大的
        \omterm{def:g7:atoms-and-molecules:atom}{原子}。
  \item 检查该步之后的电荷：给出\omterm{def:g10:lewis-and-shape:lone-pair}{孤电子对}的\omterm{def:g7:atoms-and-molecules:atom}{原子}
        正电荷增加 1，接受\omterm{def:g9:inside-the-atom:nucleus}{电子}对、形成\omterm{def:g10:lewis-and-shape:lone-pair}{孤电子对}的\omterm{def:g7:atoms-and-molecules:atom}{原子}
        负电荷增加 1；总电荷保持不变。
\end{enumerate}
\end{method}

\begin{proposition}[每一步中电荷守恒]\label{prop:g12:curly-arrows:charge}
物种的总电荷在
机理的每一步前后都相同。
\end{proposition}

\begin{proof}
\omterm{def:g12:curly-arrows:curly-arrow}{弯箭头}只移动原本就存在的\omterm{def:g9:inside-the-atom:nucleus}{电子}；
既不产生也不消灭\omterm{def:g9:inside-the-atom:nucleus}{电子}，\omterm{def:g9:inside-the-atom:nucleus}{原子核}也不变化。
\end{proof}

\begin{example}[第一个箭头]\label{ex:g12:curly-arrows:first}
氢\omterm{def:g9:ions:ion}{离子} \ce{H+} 与氢氧根\omterm{def:g9:ions:ion}{离子} \ce{OH-} 相遇。氧的一对\omterm{def:g10:lewis-and-shape:lone-pair}{孤电子对}
与没有\omterm{def:g9:inside-the-atom:nucleus}{电子}的氢\omterm{def:g9:ions:ion}{离子}形成一个键：
\ce{H+ + OH- -> H2O}。画一个箭头，从氧的\omterm{def:g10:lewis-and-shape:lone-pair}{孤电子对}指向
\ce{H+}；电荷 $+1$ 与 $-1$ 在水\omterm{def:g7:atoms-and-molecules:molecule}{分子}中变为 0。
\end{example}

\begin{omfigure}
\setchemfig{atom sep=2.1em}
\chemfig{@{hp}H^{+}}\qquad\raisebox{0.2ex}{$+$}\qquad
\chemfig{H-@{ox}\charge{90=\:,0=\:,270=\:}{O}^{-}}
\qquad\raisebox{0.2ex}{$\longrightarrow$}\qquad
\chemfig{H-\charge{90=\:,270=\:}{O}-H}
\chemmove{
  \draw[->, thick, omProp, shorten <=4pt, shorten >=3pt]
    ($(ox)+(0,0.25)$) .. controls +(110:9mm) and +(60:9mm) .. (hp);}

{\small 最简单的一步：氢氧根\omterm{def:g9:ions:ion}{离子}的一对\omterm{def:g10:lewis-and-shape:lone-pair}{孤电子对}移动，以形成
与氢\omterm{def:g9:ions:ion}{离子}之间的新 \ce{O-H} 键。}
\end{omfigure}

\section{基本步骤与中间体}

\begin{definition}[机理、基元反应、中间体]\label{def:g12:curly-arrows:mechanism}
某一反应的\emph{反应机理}\index{反应机理}是
\omterm{def:g7:chemical-reactions:reactant}{反应物}转变为生成物所经历的一系列步骤。每一个
\emph{基元反应}\index{基元反应}是一次单独的事件，其中
若干对\omterm{def:g9:inside-the-atom:nucleus}{电子}同时发生移动。在某一步中生成、又在后一步中被消耗的
物种称为\emph{反应中间体}\index{反应中间体}：
它不出现在总反应方程式中。
\end{definition}


\begin{example}[碳正离子中间体]\label{ex:g12:curly-arrows:carbocation}
溴化氢与乙烯发生\omterm{def:g12:curly-arrows:categories}{加成}时，反应分两步进行。
第一步中，乙烯的\omterm{def:g10:lewis-and-shape:multiple-bond}{双键}把一对\omterm{def:g9:inside-the-atom:nucleus}{电子}给予
\ce{HBr} 中的氢\omterm{def:g7:atoms-and-molecules:atom}{原子}，\ce{H-Br} 的\omterm{def:g9:inside-the-atom:nucleus}{电子}对归于溴\omterm{def:g7:atoms-and-molecules:atom}{原子}，溴以
\ce{Br-} 的形式离去；失去\omterm{def:g10:lewis-and-shape:multiple-bond}{双键}共用\omterm{def:g9:inside-the-atom:nucleus}{电子}对的碳\omterm{def:g7:atoms-and-molecules:atom}{原子}只剩
六个\omterm{def:g9:inside-the-atom:nucleus}{电子}并带正电荷：碳正\omterm{def:g9:ions:ion}{离子}
\ce{CH3-CH2+}。第二步中，\ce{Br-} 的一对\omterm{def:g10:lewis-and-shape:lone-pair}{孤电子对}与
该碳\omterm{def:g7:atoms-and-molecules:atom}{原子}成键。碳正\omterm{def:g9:ions:ion}{离子}是一种中间体。
\end{example}

\section{三类反应}

\begin{definition}[取代、加成、消除]\label{def:g12:curly-arrows:categories}
在\emph{取代}\index{取代}反应中，\omterm{def:g7:atoms-and-molecules:molecule}{分子}中的一个\omterm{def:g7:atoms-and-molecules:atom}{原子}或\omterm{def:g7:atoms-and-molecules:atom}{原子}团
被另一\omterm{def:g7:atoms-and-molecules:atom}{原子}或\omterm{def:g7:atoms-and-molecules:atom}{原子}团所替代。在\emph{加成}\index{加成}反应中，
两个\omterm{def:g7:atoms-and-molecules:molecule}{分子}结合为一个\omterm{def:g7:atoms-and-molecules:molecule}{分子}，其中一个多重键变为单键。在
\emph{消除}\index{消除}反应中，\omterm{def:g7:atoms-and-molecules:molecule}{分子}中脱去一个小\omterm{def:g7:atoms-and-molecules:molecule}{分子}（通常是水），
单键变为多重键。
\end{definition}

\begin{omfigure}
\setchemfig{atom sep=2em}
\small
\textbf{\omterm{def:g12:curly-arrows:categories}{取代}}（一步完成）：\ce{CH3-CH2-Br + OH- -> CH3-CH2-OH + Br-}

\medskip
\chemfig{H-@{o1}\charge{90=\:,0=\:,270=\:}{O}^{-}}\qquad
\chemfig{H_3C-@{c1}C(-[2]H)(-[6]H)-[@{b1}]@{br1}\charge{90=\:,0=\:,270=\:}{Br}}
\qquad$\longrightarrow$\qquad
\chemfig{H_3C-C(-[2]H)(-[6]H)-O-H}\quad$+$\quad\ce{Br-}
\chemmove{
  \draw[->, thick, omProp, shorten <=4pt, shorten >=3pt]
    ($(o1)+(0.15,0.2)$) .. controls +(60:9mm) and +(120:9mm) .. (c1);
  \draw[->, thick, omDef, shorten <=2pt, shorten >=4pt]
    (b1) .. controls +(-90:7mm) and +(-90:7mm) .. (br1);}

\vspace{6mm}
\textbf{\omterm{def:g12:curly-arrows:categories}{加成}}（两步完成）：\ce{CH2=CH2 + HBr -> CH3-CH2-Br}

\vspace{9mm}
\chemfig{H_2C=[@{pi}]CH_2}\quad$+$\quad
\chemfig{@{h2}H-[@{b2}]@{br2}Br}
\qquad$\longrightarrow$\qquad
\chemfig{H_3C-@{cp}\charge{90:2pt=$\scriptstyle +$}{C}H_2}\quad$+$\quad
\chemfig{@{brm}\charge{90=\:,180=\:,270=\:,0=\:}{Br}^{-}}
\qquad$\longrightarrow$\qquad \ce{CH3-CH2-Br}
\chemmove{
  \draw[->, thick, omProp, shorten <=2pt, shorten >=3pt]
    (pi) .. controls +(90:10mm) and +(110:9mm) .. (h2);
  \draw[->, thick, omDef, shorten <=2pt, shorten >=4pt]
    (b2) .. controls +(-90:7mm) and +(-90:7mm) .. (br2);
  \draw[->, thick, omProp, shorten <=4pt, shorten >=3pt]
    ($(brm)+(0,-0.25)$) .. controls +(-120:8mm) and +(-60:8mm) .. (cp);}

\vspace{6mm}
\textbf{\omterm{def:g12:curly-arrows:categories}{消除}}（乙醇的酸催化脱水，三步完成）：
\ce{CH3-CH2-OH -> CH2=CH2 + H2O}

\medskip
\begin{tabular}{@{}l@{}}
1.\ 氧\omterm{def:g7:atoms-and-molecules:atom}{原子}获得一个氢\omterm{def:g9:ions:ion}{离子}：\ce{CH3-CH2-OH + H+ -> CH3-CH2-OH2+}；\\
2.\ \ce{C-O} \omterm{def:g9:inside-the-atom:nucleus}{电子}对与水一起离去：\ce{CH3-CH2-OH2+ -> CH3-CH2+ + H2O}；\\
3.\ 相邻的 \ce{C-H} \omterm{def:g9:inside-the-atom:nucleus}{电子}对形成\omterm{def:g10:lewis-and-shape:multiple-bond}{双键}，同时
    \ce{H+} 离去：\\
\phantom{3.\ }\ce{CH3-CH2+ -> CH2=CH2 + H+}。
\end{tabular}

\medskip
{\small 三种机理。橙色箭头：给\omterm{def:g9:inside-the-atom:nucleus}{电子}对形成新键。
蓝色箭头：键断裂，其\omterm{def:g9:inside-the-atom:nucleus}{电子}对归于\omterm{def:g11:polarity-and-cohesion:electronegativity}{电负性}更大的
\omterm{def:g7:atoms-and-molecules:atom}{原子}。\omterm{def:g12:curly-arrows:categories}{消除}反应中，氢\omterm{def:g9:ions:ion}{离子}最后被重新释放。}
\end{omfigure}


\begin{remark}[改变骨架或官能团]\label{rem:g12:curly-arrows:skeleton}
有些反应只改变\omterm{def:g11:functional-groups:functional-group}{官能团}，而保持碳骨架不变：
溴乙烷的\omterm{def:g12:curly-arrows:categories}{取代}反应将\omterm{def:g11:functional-groups:halogenoalkane}{卤代烷}转化为醇，
反应前后各含两个碳\omterm{def:g7:atoms-and-molecules:atom}{原子}。另一些反应则构建或断裂
骨架：形成新的碳碳键会使碳链加长，这是
由小\omterm{def:g7:atoms-and-molecules:molecule}{分子}\omterm{def:g10:synthesis-yield:synthesis}{合成}大\omterm{def:g7:atoms-and-molecules:molecule}{分子}的关键步骤。设计\omterm{def:g10:synthesis-yield:synthesis}{合成}路线
需要按正确顺序选择这两类反应。
\end{remark}

\section{练习}

\begin{exercise}[$\star$]\label{exo:g12:curly-arrows:1}
在下列每种物种中，指出一个\omterm{def:g12:curly-arrows:site}{给电子位点}，若存在，再指出一个
\omterm{def:g12:curly-arrows:site}{受电子位点}：\ce{OH-}；\ce{H+}；\ce{CH2=CH2}；\ce{CH3-Cl}；\ce{H2O}。
\end{exercise}

\begin{exercise}[$\star$]\label{exo:g12:curly-arrows:2}
下列反应分别属于\omterm{def:g12:curly-arrows:categories}{取代}、\omterm{def:g12:curly-arrows:categories}{加成}还是\omterm{def:g12:curly-arrows:categories}{消除}反应？
(a) \ce{CH3-CH2-Cl + OH- -> CH3-CH2-OH + Cl-}；
(b) \ce{CH2=CH2 + H2O -> CH3-CH2-OH}；
(c) \ce{CH3-CH2-OH -> CH2=CH2 + H2O}；
(d) \ce{CH2=CH2 + Br2 -> CH2Br-CH2Br}。
\end{exercise}

\begin{exercise}[$\star$]\label{exo:g12:curly-arrows:3}
\omterm{def:g12:curly-arrows:curly-arrow}{弯箭头}从哪里开始，在哪里结束？
从氧\omterm{def:g7:atoms-and-molecules:atom}{原子}的\omterm{def:g10:lewis-and-shape:lone-pair}{孤电子对}指向碳\omterm{def:g7:atoms-and-molecules:atom}{原子}的\omterm{def:g12:curly-arrows:curly-arrow}{弯箭头}表示什么？
\end{exercise}

\begin{exercise}[$\star$]\label{exo:g12:curly-arrows:4}
在反应 \ce{H+ + NH3 -> NH4+} 中，哪种物种提供\omterm{def:g9:inside-the-atom:nucleus}{电子}对？
画出\omterm{def:g12:curly-arrows:curly-arrow}{弯箭头}。生成物带何种电荷？
\end{exercise}

\begin{exercise}[$\star$]\label{exo:g12:curly-arrows:5}
什么是\omterm{def:g12:curly-arrows:mechanism}{反应中间体}？请写出
\ce{HBr} 与乙烯\omterm{def:g12:curly-arrows:categories}{加成}反应的中间体。
\end{exercise}

\begin{exercise}[$\star\star$]\label{exo:g12:curly-arrows:6}
画出\omterm{def:g12:curly-arrows:categories}{取代}反应
\ce{CH3-Br + OH- -> CH3-OH + Br-}（一步完成）的\omterm{def:g12:curly-arrows:curly-arrow}{弯箭头}，并检查电荷。
\end{exercise}

\begin{exercise}[$\star\star$]\label{exo:g12:curly-arrows:7}
在 \ce{HBr} 与乙烯\omterm{def:g12:curly-arrows:categories}{加成}的第一步中，为什么
\ce{H-Br} \omterm{def:g9:inside-the-atom:nucleus}{电子}对归属于溴而不是氢？
该步每种产物分别带何种电荷？
\end{exercise}

\begin{exercise}[$\star\star$]\label{exo:g12:curly-arrows:8}
写出乙醇脱水的三步反应，并对每一步
指出哪个物种是\omterm{def:g12:curly-arrows:site}{给电子位点}、哪个是\omterm{def:g12:curly-arrows:site}{受电子位点}。
为什么氢\omterm{def:g9:ions:ion}{离子}不写入总反应方程式？
\end{exercise}

\begin{exercise}[$\star\star$]\label{exo:g12:curly-arrows:9}
步骤 \ce{CH3-CH2+ + H2O -> CH3-CH2-OH2+} 紧接在
乙烯于水中\omterm{def:g12:curly-arrows:categories}{加成}氢\omterm{def:g9:ions:ion}{离子}之后。画出\omterm{def:g12:curly-arrows:curly-arrow}{弯箭头}。
再失去 \ce{H+} 后生成什么？写出总反应方程式，并
说明其类别。
\end{exercise}

\begin{exercise}[$\star\star$]\label{exo:g12:curly-arrows:10}
阅读溴乙烷\omterm{def:g12:curly-arrows:categories}{取代}反应的机理：有几对
\omterm{def:g9:inside-the-atom:nucleus}{电子}对发生移动？哪个键形成，哪个键断裂？是否有中间体？
\end{exercise}

\begin{exercise}[$\star\star$]\label{exo:g12:curly-arrows:11}
% ledger: en:C, en:Cl, en:Br
氯甲烷与氢氧根\omterm{def:g9:ions:ion}{离子}的反应比溴甲烷
慢。利用\omterm{def:g11:polarity-and-cohesion:electronegativity}{电负性}（C 2.55，Cl 3.16，Br 2.96），
解释哪个碳\omterm{def:g7:atoms-and-molecules:atom}{原子}的 $\delta^+$ 更大。这能否
解释二者速率的差异？（还应考虑离去\omterm{def:g9:ions:ion}{离子}夺取\omterm{def:g9:inside-the-atom:nucleus}{电子}对的难易程度。）
\end{exercise}

\begin{exercise}[$\star\star\star$]\label{exo:g12:curly-arrows:12}
氯化氢\omterm{def:g12:curly-arrows:categories}{加成}到丙烯 \ce{CH2=CH-CH3} 上。第一步中
氢\omterm{def:g9:ions:ion}{离子}可与\omterm{def:g10:lewis-and-shape:multiple-bond}{双键}两端的任一碳\omterm{def:g7:atoms-and-molecules:atom}{原子}成键，
生成两种可能的碳正\omterm{def:g9:ions:ion}{离子}。画出这两种碳正\omterm{def:g9:ions:ion}{离子}及两种可能的产物，并
为它们命名。（哪一种生成最多，在本科一年级分册中学习。）
\end{exercise}

\begin{exercise}[$\star\star\star$]\label{exo:g12:curly-arrows:13}
一名学生画出溴乙烷\omterm{def:g12:curly-arrows:categories}{取代}反应的第一步，
其\omterm{def:g12:curly-arrows:curly-arrow}{弯箭头}从溴\omterm{def:g7:atoms-and-molecules:atom}{原子}指向 \ce{OH-} 中的氧\omterm{def:g7:atoms-and-molecules:atom}{原子}。请解释
其中的两处错误。
\end{exercise}

\begin{exercise}[$\star\star\star$]\label{exo:g12:curly-arrows:14}
由\omterm{def:g11:functional-groups:carboxylic-acid}{羧酸}与醇生成酯的过程中，醇的氧\omterm{def:g7:atoms-and-molecules:atom}{原子}
进攻 \ce{C=O} 基团中的碳\omterm{def:g7:atoms-and-molecules:atom}{原子}。利用
\omterm{def:g11:polarity-and-cohesion:polar-bond}{部分电荷}，解释为什么该碳\omterm{def:g7:atoms-and-molecules:atom}{原子}是\omterm{def:g12:curly-arrows:site}{受电子位点}、
该氧\omterm{def:g7:atoms-and-molecules:atom}{原子}是\omterm{def:g12:curly-arrows:site}{给电子位点}。总反应失去的是哪种小\omterm{def:g7:atoms-and-molecules:molecule}{分子}，以及
整个反应属于哪一类？
\end{exercise}

\begin{exercise}[$\star\star\star$]\label{exo:g12:curly-arrows:15}
乙烯与溴 \ce{Br2}（\omterm{def:g11:polarity-and-cohesion:polar-molecule}{非极性分子}）反应。当 \ce{Br2}
\omterm{def:g7:atoms-and-molecules:molecule}{分子}接近\omterm{def:g10:lewis-and-shape:multiple-bond}{双键}时，\ce{Br-Br} 的\omterm{def:g9:inside-the-atom:nucleus}{电子}被推离靠近的溴\omterm{def:g7:atoms-and-molecules:atom}{原子}。
解释为何这会形成一个
\omterm{def:g12:curly-arrows:site}{受电子位点}，并提出机理的第一步。
\end{exercise}

\section{问题：制备溴乙烷}

\begin{problem}[{周末问题——从 10 克乙烯出发：能
得到多少溴乙烷？}]
\label{pb:g12:curly-arrows:1}
% ledger: aw:C, aw:H, aw:Br, en:H, en:Br, en:C
溴乙烷是许多\omterm{def:g10:synthesis-yield:synthesis}{合成}的起始原料，可通过
溴化氢与乙烯\omterm{def:g12:curly-arrows:categories}{加成}制得：\ce{CH2=CH2 + HBr -> CH3-CH2-Br}。一位化学家
用 \qty{10.0}{g} 乙烯和过量的溴化氢进行反应，
其\omterm{def:g10:synthesis-yield:yield}{产率}为 \qty{75}{\percent}。

\medskip
\noindent\textbf{第一部分——位点。}
\begin{enumerate}
  \item 画出乙烯和溴化氢的\omterm{def:g10:lewis-and-shape:lewis-structure}{路易斯结构}。
  \item 乙烯的\omterm{def:g12:curly-arrows:site}{给电子位点}在哪里？为什么？
  \item 计算 \ce{H-Br} 键的\omterm{def:g11:polarity-and-cohesion:electronegativity}{电负性}之差
        （H 2.20，Br 2.96）。哪个\omterm{def:g7:atoms-and-molecules:atom}{原子}带 $\delta^+$ 电荷？
  \item \ce{HBr} 中哪个\omterm{def:g7:atoms-and-molecules:atom}{原子}是\omterm{def:g12:curly-arrows:site}{受电子位点}？
\end{enumerate}

\medskip
\noindent\textbf{第二部分——\omterm{def:g12:curly-arrows:mechanism}{反应机理}。}
\begin{enumerate}[resume]
  \item 用两个\omterm{def:g12:curly-arrows:curly-arrow}{弯箭头}画出第一步反应。
  \item 写出生成的两种物种及其电荷。
  \item 检验第一步中电荷是否守恒。
  \item 用\omterm{def:g12:curly-arrows:curly-arrow}{弯箭头}画出第二步反应。
  \item 中间体中带正电的碳\omterm{def:g7:atoms-and-molecules:atom}{原子}周围有多少个\omterm{def:g9:inside-the-atom:nucleus}{电子}？
        这种中间体为何如此活泼？
\end{enumerate}

\medskip
\noindent\textbf{第三部分——属于哪类反应？}
\begin{enumerate}[resume]
  \item 该\omterm{def:g12:curly-arrows:mechanism}{反应机理}中的中间体是什么？
  \item 它为什么不出现在总反应方程式中？
  \item 该反应属于哪一类？
  \item 该反应改变了碳骨架，还是只改变
        \omterm{def:g11:functional-groups:functional-group}{官能团}？
\end{enumerate}

\medskip
\noindent\textbf{第四部分——能得到多少产物？}
\begin{enumerate}[resume]
  \item 计算乙烯和溴乙烷的\omterm{def:g10:the-mole:molar-mass}{摩尔质量}。
  \item 计算乙烯的\omterm{def:g10:the-mole:amount}{物质的量}。
  \item 填写\omterm{def:g10:reaction-progress-table:extent}{进度表}，哪种\omterm{def:g7:chemical-reactions:reactant}{反应物}是\omterm{def:g10:reaction-progress-table:limiting}{限量反应物}？
  \item 计算溴乙烷的最大质量。
  \item 计算\omterm{def:g10:synthesis-yield:yield}{产率}为 \qty{75}{\percent} 时得到的质量。
  \item 给出最终答案：化学家能得到
        多少质量的溴乙烷？
\end{enumerate}
\end{problem}
